\section{Introduction}
\label{sec:introduction}

Digital laboratories join instrument control, physical-resource models, workflow orchestration, sample and event records, and increasingly AI-assisted method generation. Standards such as SiLA and SiLA~2 expose device capabilities through shared communication models, and open frameworks such as PyLabRobot provide hardware-agnostic programming abstractions \cite{bar2012sila,juchli2022sila2,porr2020gateway,bromig2022manager,wierenga2023pylabrobot}. Autonomous-laboratory architectures extend these foundations through adaptive planning, automated experimentation, and generalized scheduling \cite{abolhasani2023rise,hung2024community,canty2023autonomy,fei2024alabos,cousty2024glas}. These systems depend on digital objects whose operational meaning remains coupled to physical state. An API declaration must correspond to executable access. A method artifact must identify its deployment. A resource record must encode valid geometry. A trace must support reconstruction of material state. A simulation result must remain bounded to the stage it evaluated.

The literature principally presents standards, architectures, software platforms, and autonomous demonstrations. Public practitioner discussions preserve another form of digital-laboratory evidence. They record the file structures, version transitions, licence conditions, geometry defects, state divergences, recovery procedures, and validation disputes that appear during integration and maintenance. These discussions also record successful transformations of local diagnosis into public code, documentation, examples, and post-mortems. Their scientific value lies in the specificity of the represented objects and relationships.

Forum activity does not provide the operational denominators required for prevalence, reliability, or market comparisons. Products differ in installed base, public participation, complexity, and reporting propensity. Discussions can become publicly censored after work moves into private support channels or local implementation. The appropriate analytical strategy is therefore to preserve explicit units, source provenance, uncertainty, validation stage, and prohibited inference. This study treats public discussions as inputs to a reproducible evidence infrastructure, not as a census of laboratory failures.

The paper addresses three questions.
\begin{enumerate}[label=\textbf{RQ\arabic*.}]
  \item Which recurrent digital and physical bottlenecks can be represented consistently across public laboratory-automation discussions?
  \item Which bounded metrics can be computed from machine-readable case records without converting discussion recurrence into operational incidence?
  \item Which reusable data structures and software controls can support prospective digital-laboratory research and community knowledge stewardship?
\end{enumerate}

The contribution combines empirical analysis with an executable research resource. The release contains a ten-construct taxonomy, thread- and episode-level derived records, field-level metric inputs, robustness analyses, source and quotation audits, a publication claim ledger, deterministic table and figure generation, regression tests, JSON Schemas, and governed example records. The central empirical mechanism links execution observability with recovery after irreversible physical actions. The resource contribution defines how evidence about interfaces, deployments, physical resources, events, validation stages, and resolved knowledge can remain auditable across analysis, publication, and later correction.

% claim: C11
The resulting model is boundary-centered. Digital-laboratory reliability depends on transitions from public knowledge to maintained artifacts, listed interfaces to testable access, source to identified deployment, resource values to validated physical models, telemetry to adjudicable evidence, interrupted execution to verified disposition, and generated methods to experimentally supported claims. These transitions provide a common representation for otherwise heterogeneous operational problems and a basis for prospective measurement with real run and incident denominators.
